\section{Introduction}
\label{sec:intro}

Semantic correspondence is the task of identifying pixel-level correspondences between semantically equivalent parts of objects across different
images. Given a keypoint on an object in a source image, the objective is to localize the corresponding semantic part in a target image,
even when the two images differ significantly in appearance, pose, scale, or domain. This problem is fundamental to many computer vision applications,
including object recognition, image alignment, 3D reconstruction, and visual reasoning.

Unlike classical geometric correspondence, semantic correspondence must cope with substantial intra-class variability and non-rigid deformations.
Objects may appear under different viewpoints or exhibit strong shape variations, and semantically similar parts may not be geometrically aligned.
In addition, models must correctly distinguish between symmetric or visually similar parts, such as left and right limbs, which further increases
the ambiguity of the task.

Recent work has shown that semantic correspondence can emerge naturally from large pretrained visual models~\cite{tang2023emergent, zhang2023tale, zhang2024telling}.
Vision Transformers trained with self-supervised objectives, such as DINO~\cite{caron2021emerging} and its successors~\cite{oquab2023dinov2, simeoni2025dinov3},
learn feature representations that capture high-level semantic structure without requiring explicit correspondence annotations.
Similarly, foundation models such as Segment Anything~\cite{kirillov2023segment} provide powerful dense representations that generalize across tasks
and domains. These observations suggest that pretrained foundation models may serve as strong feature extractors for semantic correspondence,
even in training-free or weakly supervised settings.

In this work, we perform a systematic evaluation of semantic correspondence using pretrained visual foundation models. Using the SPair-71k benchmark~\cite{min2019spair},
we first establish a training-free baseline by extracting frozen dense features and performing correspondence via cosine similarity and global argmax
matching, following prior work~\cite{tang2023emergent}. We then investigate the impact of light fine-tuning by unfreezing and adapting the final layers
of the backbone using contrastive objective and keypoint supervision. Finally, we study the effect of replacing the standard argmax prediction with
a windowed soft-argmax strategy~\cite{zhang2024telling}, which enables sub-pixel refinement and improves robustness to noisy similarity maps.

Our experimental analysis aims to answer three key questions: (i) how well different foundation models support semantic correspondence without
fine-tuning, (ii) how much performance can be gained through limited adaptation of pretrained representations, and (iii) how does the prediction rule
influence correspondence accuracy. Our results show that DINO-based models significantly outperform SAM for semantic correspondence, with fine-tuned
DINOv2 and DINOv3 achieving comparable performance at PCK@0.1 (74.89\% and 75.77\%, respectively), while DINOv2 attains higher accuracy at stricter
thresholds. Moreover, light fine-tuning and windowed soft-argmax consistently improve performance across models. These findings highlight the
effectiveness of foundation-model-based approaches and provide practical insights into the design of semantic correspondence pipelines.

